% Full derivation figure for the worked example (numbers from
% code/ch12_velocity_obstacles.py): (a) workspace, (b) velocity space of A.
\begin{tikzpicture}
% ---------------- (a) workspace, 0.55 cm per metre, A at the origin ----------
\begin{scope}[scale=0.55]
  \coordinate (A) at (0,0);
  \coordinate (B) at (5,-5);
  % collision cone at A: legs (0.8,-0.6) and (0.6,-0.8), tangency at distance 7
  \fill[sbOrange!15] (A) -- (7.6,-5.7) -- (5.7,-7.6) -- cycle;
  \draw[sbOrange,thick] (A) -- (7.6,-5.7);
  \draw[sbOrange,thick] (A) -- (5.7,-7.6);
  \fill[sbOrange] (5.6,-4.2) circle (1.6pt);
  \fill[sbOrange] (4.2,-5.6) circle (1.6pt);
  \draw[dashed,sbGray] (A) -- (B);
  \draw[sbOrange,dashed,thick] (B) circle (1);
  \draw[draw=sbAgentB,fill=sbAgentB!35,thick] (B) circle (0.5);
  \draw[draw=sbAgentA,fill=sbAgentA!35,thick] (A) circle (0.5);
  \node[font=\footnotesize] at (A) {$A$};
  \node[font=\footnotesize] at (B) {$B$};
  % velocity arrows (drawn at twice the metric scale)
  \draw[sbvecA] (A) -- ++(2,0) node[pos=1,right,sbannot] {$\vel_A$};
  \draw[sbvecB] (B) -- ++(0,1.8) node[pos=1,right,sbannot] {$\vel_B$};
  \draw[sbGray] (B) -- ++(20:1) node[pos=1.35,sbannot] {$R$};
  \node[sbannot] at (1.6,-2.9) {$d$};
  \draw[sbGray] (-45:2.6) arc[start angle=-45,end angle=-36.87,radius=2.6];
  \node[sbannot] at (2.75,-1.75) {$\theta$};
  \node[sbannot,text=sbOrange!80!black] at (3.4,-6.9) {$CC_{A|B}$};
  \node[font=\footnotesize] at (3.0,-9.2) {(a) workspace};
\end{scope}
% ---------------- (b) velocity space of A, 2.6 cm per m/s ---------------------
\begin{scope}[xshift=7.6cm,yshift=-3.55cm,scale=2.6]
  \clip (-0.42,-1.05) rectangle (1.85,1.3);
  \draw[->,black!45] (-0.38,0) -- (1.8,0) node[pos=1,below,sbannot] {$v_x$};
  \draw[->,black!45] (0,-1.0) -- (0,1.25) node[pos=1,left,sbannot] {$v_y$};
  \draw[black!45,dotted,thick] (0,0) circle (1.5);
  \coordinate (P) at (0,0.9);
  % untruncated wedge: legs (0.8,-0.6), (0.6,-0.8) from the apex
  \fill[sbOrange!15] (P) -- (2.4,-0.9) -- (1.8,-1.5) -- cycle;
  \draw[sbOrange,thin] (P) -- (2.4,-0.9);
  \draw[sbOrange,thin] (P) -- (1.8,-1.5);
  % truncated region: legs beyond the tangency points (1.12,0.06), (0.84,-0.22)
  % and the near arc of the disc centre (1,-0.1), radius 0.2
  \fill[sbOrange!45,draw=sbOrange,thick] (2.4,-0.9) -- (1.12,0.06)
     arc[start angle=53.13,end angle=216.87,radius=0.2] -- (1.8,-1.5) -- cycle;
  \draw[sbOrange,thin] (1,-0.1) circle (0.2);
  \draw[dashed,sbGray] (P) -- (1.7,-0.8);
  \draw[sbGray,thin] (P) -- (1,0);
  \node[sbannot] at (0.42,0.55) {$\vel_A^{\mathrm{pref}}-\vel_B$};
  \fill[sbOrange] (P) circle (1.4pt);
  \node[sbannot,left] at (P) {$\vel_B$};
  \fill[sbRed] (1,0) circle (1.4pt);
  \node[sbannot,text=sbRed,below right] at (1.0,-0.005) {$\vel_A^{\mathrm{pref}}$};
  \fill[sbGreen] (1,0.1) circle (1.4pt);
  \node[sbannot,text=sbGreen,above right] at (1.0,0.11) {$\vel^*$};
  \fill[sbBlue] (0.897,0.078) circle (1.4pt);
  \node[sbannot,text=sbBlue,above left] at (0.89,0.09) {$\vel_A'$};
  \node[sbannot,text=sbOrange!80!black] at (1.55,-0.25) {$\VO^{\ttc}_{A|B}$};
  \node[sbannot,text=sbOrange!80!black] at (0.58,0.36) {$\VO_{A|B}$};
  \node[sbannot] at (0.45,-0.25) {$\theta$};
  \node[sbannot] at (1.15,1.1) {$\norm{\vel}=v_{\max}$};
  \node[font=\footnotesize] at (0.75,-0.95) {(b) velocity space of $A$};
\end{scope}
\end{tikzpicture}
